% Fiche SystemVerilog pour la vérification + UVM — recto-verso A4 paysage
% Sources : AED_SV_UVM_Cheatsheet (pages 3-6), cours Phelma B (classes, aléatoire),
%   C (UVM, séquences), D (couverture, scoreboard), code des TP 1 à 4.
% Complète la fiche Verilog (../Verilog) : le langage, la conception, les interfaces,
%   les threads et les tableaux dynamiques sont là-bas ; ici seulement la vérification.
% Compiler : make (ici ou à la racine), ou xelatex uvm_cheat.tex — une seule passe suffit
% Apparence : ../librairie (fiche.sty + theme-uvm.tex = thème Verilog + catégories
%   oop rnd cov sva uvm obs). Macros : colonne[titre][couleur] · bloc{titre}{cat} · code
%   · \V|...| · «...» = à remplir · \sub · \note · \piege · schema
%   optab étiquette | texte · tab3 étiquette | code | remarque (largeurs mesurées)
% Le code est en chasse fixe : commentaires alignés à la main (les « » ne comptent pas).
% Après compilation : grep "\[fiche\]" uvm_cheat.log (place libre, collisions).
\documentclass{article}
\usepackage[uvm]{../librairie/fiche}   % apparence : ../librairie (fiche.sty + theme-uvm.tex)
% --- Langage SystemVerilog + UVM pour listings ---
% Classes de mots-clés : [1] rose · [2] cyan · [3] vert · [4] couleur du thème
% (styles dans ../librairie/fiche.sty). Chargé par uvm_cheat.tex.
\lstdefinelanguage{CODE}{
  sensitive=true, alsoletter={\$`},
  morekeywords=[1]{module,endmodule,input,output,interface,endinterface,package,
    endpackage,import,class,endclass,extends,virtual,function,endfunction,task,
    endtask,new,super,this,local,protected,static,extern,rand,randc,constraint,
    inside,dist,solve,before,soft,with,foreach,if,else,begin,end,return,forever,
    repeat,fork,join,join_any,join_none,initial,always,typedef,enum,posedge,
    negedge,null,type,randcase,case,endcase,for,while,do,covergroup,endgroup,
    coverpoint,bins,ignore_bins,illegal_bins,cross,binsof,intersect,option,
    property,endproperty,sequence,endsequence,assert,assume,cover,disable,iff,
    bind,default,clocking,endclocking,const},
  morekeywords=[2]{bit,logic,int,byte,string,event,real,integer,reg,wire,
    unsigned,signed,void,uvm_object,uvm_component,uvm_test,uvm_env,uvm_agent,
    uvm_driver,uvm_monitor,uvm_sequencer,uvm_sequence,uvm_sequence_item,
    uvm_scoreboard,uvm_phase,uvm_analysis_port,uvm_config_db,uvm_object_wrapper},
  morekeywords=[3]{\$display,\$sformatf,\$cast,\$error,\$rose,\$fell,\$stable,
    \$past,\$onehot,\$onehot0,\$isunknown,\$urandom,\$urandom_range,\$countones,
    `uvm_component_utils,`uvm_object_utils,`uvm_object_utils_begin,
    `uvm_object_utils_end,`uvm_field_int,`uvm_field_enum,`uvm_info,`uvm_warning,
    `uvm_error,`uvm_fatal,`uvm_do,`uvm_do_with,`uvm_analysis_imp_decl,`include,
    `timescale,run_test},
  morecomment=[l]{//}, morecomment=[s]{/*}{*/}, morestring=[b]",
}
   % coloration syntaxique : \lstdefinelanguage{CODE}

\begin{document}
%%%%%%%%%%%%%%%%%%%%%%%%%%%%%%%%%%%%%%%%%%%%%%%%%%%%%%%%%%%%%%%%%%%%%%%%%%%%%%
%                RECTO : SYSTEMVERILOG POUR LA VÉRIFICATION
%%%%%%%%%%%%%%%%%%%%%%%%%%%%%%%%%%%%%%%%%%%%%%%%%%%%%%%%%%%%%%%%%%%%%%%%%%%%%%
\entete{SystemVerilog · vérification}{recto · classes, aléatoire, couverture, assertions}
\noindent
% ======================= colonne 1 : classes =======================
\begin{colonne}[classes][col-a]
\begin{bloc}{Classes et objets}{oop}
\begin{code}
class Packet;
  bit [31:0] addr;             // public (défaut)
  local     int id;            // cette classe seule
  protected int crc;           // + classes filles
  static    int count = 0;     // partagé par tous
  function new(bit [31:0] a = 0);  // constructeur,
    addr = a;  count++;            // sans type retour
  endfunction
  virtual function void print();   // redéfinissable
    $display("addr=%h", addr);
  endfunction
  extern function int calc();  // corps hors classe
endclass
function int Packet::calc();   // corps de calc()
  return addr ^ crc;
endfunction
\end{code}
\sub{Héritage}
\begin{code}
class BigPacket extends Packet;
  rand bit [7:0] len;
  function new(bit [31:0] a = 0);
    super.new(a);              // constructeur parent
  endfunction
  virtual function void print();
    super.print();             // méthode du parent
    $display("len=%0d", len);
  endfunction
endclass
\end{code}
\sub{Handles, polymorphisme, cast}
\begin{code}
Packet    p;                   // handle = null
BigPacket b = new(32'h10);     // alloue l'objet
p = b;                         // fille -> parent : ok
p.print();                     // BigPacket::print
if (!$cast(b, p))              // parent -> fille :
  $error("mauvais type");      //   vérifié au run
n = Packet::count;             // membre statique
\end{code}
\note{Pas de destructeur (ramasse-miettes) · \V|this| objet courant ·
\V|super| parent · \V|typedef class C;| déclaration anticipée.}
\sub{Classe paramétrée}
\begin{code}
class Stack #(type T = int, int N = 8);
  T items[$:N-1];              // file bornée
  function void push(T v); items.push_back(v); endfunction
endclass
Stack #(string) s = new;
\end{code}
\end{bloc}

\begin{bloc}{Métriques de couverture}{cov}
\begin{optab}[15mm]
Code        & lignes / blocs, branches, conditions, toggle (0$\to$1, 1$\to$0),
              états et transitions des FSM. Automatique, dit ce qui a été \emph{exécuté} \\
Fonct.      & \V|covergroup| écrits à partir du plan de vérif : dit quelles
              \emph{fonctions} ont été exercées \\
Assertions  & \V|cover property| : scénarios atteints
\end{optab}
\begin{tab3}[18mm]{21mm}
Code 100\,\%   & Fonct.~<~100\,\% & tests incomplets \\
Code < 100\,\% & Fonct.~100\,\%   & plan de vérif incomplet \\
Code 100\,\%   & Fonct.~100\,\%   & bien, mais rien d'oublié ?
\end{tab3}
\note{Exécuté $\neq$ vérifié : il faut aussi des checkers (assertions, scoreboard).
\V|covergroup| : colonne 3.}
\end{bloc}
\end{colonne}\colbreak
% ======================= colonne 2 : aléatoire =======================
\begin{colonne}[aléatoire][col-b]
\begin{bloc}{rand et contraintes}{rnd}
\begin{code}
typedef enum {LOW, MID, HIGH} kind_t;
class Trans;
  rand  bit [15:0] addr;
  rand  bit [31:0] data;
  rand  kind_t     kind;
  randc bit [3:0]  id;         // cyclique (permutation)
  rand  int        q[];        // taille aussi aléatoire
  constraint c_align { addr[1:0] == 2'b00; }
  constraint c_range {
    kind == LOW  -> addr inside {[0:15]};
    kind == HIGH -> addr inside {[128:255]};
  }
  constraint c_q { q.size() inside {[1:8]};
                   foreach (q[i]) q[i] < 100; }
endclass
\end{code}
\sub{Formes de contraintes}
\begin{tab3}[17mm]{33mm}
Relation    & \V|a < b; a != 0;|          & expression booléenne \\
Ensemble    & \V|x inside {[0:9], 12};|   & intervalles et valeurs \\
Exclusion   & \V|!(x inside {[0:9]});|    & hors de l'ensemble \\
Implication & \V|c -> x == 0;|            & si c alors... \\
Condition   & \V|if (c) {..} else {..}| & idem, avec else \\
Douce       & \V|soft len < 20;|          & abandonnée si conflit \\
Ordre       & \V|solve k before a;|       & tire k d'abord \\
Tableau     & \V|foreach (q[i]) {...}|    & par élément
\end{tab3}
\sub{Distribution : dist}
\begin{code}
x dist { [100:102] := 1,     // chaque valeur : poids 1
         200       := 2,     // poids 2
         [300:304] :/ 1 };   // poids 1 partagé sur 5
// total 3+2+1 = 6 : P(200) = 2/6, P(300..304) = 1/6
\end{code}
\end{bloc}

\begin{bloc}{Tirage}{rnd}
\begin{code}
Trans t = new;
if (!t.randomize())            // 0 si les contraintes
  $error("échec");             //   sont insolubles
assert (t.randomize() with {   // contraintes en ligne
  addr < 'h100;
  kind == local::kind;         // kind de l'appelant
});
void'(t.randomize());          // ignore le résultat
t.c_align.constraint_mode(0);  // coupe une contrainte
t.data.rand_mode(0);           // data n'est plus tirée
t.randomize(addr);             // ne tire que addr
\end{code}
\begin{optab}[24mm]
\V|pre_randomize()|  & appelée juste avant le tirage (allouer un objet \V|rand|) \\
\V|post_randomize()| & juste après (calculer les champs dérivés)
\end{optab}
\begin{tab3}[0mm]{42mm}
 & \V|$urandom()|               & 32 bits non signé \\
 & \V|$urandom_range(max, min)| & entre min et max \\
 & \V|std::randomize(x) with {...}| & hors classe
\end{tab3}
\sub{randcase : actions pondérées}
\begin{code}
randcase
  3          : x = 1;          // 3/8 si en = 1, 3/4 sinon
  1          : x = 2;          // 1/8 si en = 1, 1/4 sinon
  4*en       : x = 3;          // poids = expression
endcase
\end{code}
\note{La graine change les tirages : \V|vsim -sv_seed random|.}
\end{bloc}
\end{colonne}\colbreak
% ======================= colonne 3 : couverture, assertions =======================
\begin{colonne}[couverture · SVA][col-c]
\begin{bloc}{Couverture fonctionnelle}{cov}
\begin{code}
covergroup cg_tr @(posedge clk);  // ou sample()
  cp_addr : coverpoint tr.addr {
    bins zero    = {0};
    bins low     = {[1:63], 65};  // 1 bin
    bins mid[4]  = {[64:127]};    // 4 bins égaux
    bins val[]   = {200,201,202}; // 1 bin / valeur
    ignore_bins  ign = {[250:255]};
    illegal_bins bad = {128};     // erreur si vu
  }
  cp_dir : coverpoint tr.dir;     // 1 bin / enum
  cx     : cross cp_addr, cp_dir;
endgroup
cg_tr cg = new;                   // instancier !
cg.sample();                      // échantillonner
cg.get_coverage();                // % atteint
\end{code}
\sub{Transitions et croisements}
\begin{code}
coverpoint state {
  bins t1   = (IDLE => RUN => DONE); // séquence
  bins t3   = (RUN [*3]);            // 3 fois de suite
  bins ts[] = (IDLE, RUN => DONE);   // 2 transitions
}
c : cross a, b {
  bins c1 = binsof(a.low) && binsof(b.zero);
  ignore_bins c2 = binsof(a) intersect {[100:200]};
}
\end{code}
\sub{sample() avec arguments (dans un monitor)}
\begin{code}
covergroup cg_par with function sample(bit p, byte d);
  cp_p : coverpoint p;
  cp_d : coverpoint ^d;           // parité des données
  cross cp_p, cp_d;
endgroup
cg_par cgp = new;  cgp.sample(cfg.parity, fr.data);
\end{code}
\piege{échantillonner ce qui est \emph{observé} (monitor).}
\end{bloc}

\begin{bloc}{Assertions (SVA)}{sva}
\begin{code}
// immédiate : vérifiée quand elle est exécutée
assert (a == b) else $error("a=%0d b=%0d", a, b);
// concurrente : vérifiée à chaque front
property p_req_gnt;
  @(posedge clk) disable iff (!rst_n)
  req |-> ##[1:3] gnt;            // gnt 1 à 3 cycles
endproperty
A1: assert property (p_req_gnt) else $error("gnt");
C1: cover  property (p_req_gnt);  // compte les succès
assume property (@(posedge clk) !(push && pop));
sequence s_wr; req ##1 we ##1 !req; endsequence
bind fifo fifo_chk chk_i (.*);    // greffe dans le DUT
\end{code}
\begin{tab3}[0mm]{33mm}
 & \V/a |-> b/ & b au même cycle que a \\
 & \V/a |=> b/ & b au cycle suivant (= \V/|-> ##1/) \\
 & \V|##n  ##[m:n]| & délai de n (m à n) cycles \\
 & \V|a[*n]  a[*m:n]| & a répété n fois consécutives \\
 & \V|a[->n]| & n-ième occurrence de a \\
 & \V|$rose(a) $fell(a)| & front montant / descendant \\
 & \V|$stable(a) $past(a,n)| & inchangé · valeur n cycles avant \\
 & \V|$onehot(v) $onehot0(v)| & exactement un 1 · au plus un 1 \\
 & \V|$isunknown(v)| & v contient x ou z
\end{tab3}
\note{\V|assume| : hypothèse sur les entrées (preuve formelle, Jasper).}
\end{bloc}
\end{colonne}

%%%%%%%%%%%%%%%%%%%%%%%%%%%%%%%%%%%%%%%%%%%%%%%%%%%%%%%%%%%%%%%%%%%%%%%%%%%%%%
%                            VERSO : UVM
%%%%%%%%%%%%%%%%%%%%%%%%%%%%%%%%%%%%%%%%%%%%%%%%%%%%%%%%%%%%%%%%%%%%%%%%%%%%%%
\newpage
\entete{UVM}{verso · architecture, séquences, composants}
\noindent
% ======================= colonne 4 : architecture =======================
\begin{colonne}[architecture][col-a]
\begin{bloc}{Architecture d'un banc UVM}{uvm}
\begin{schema}[x=1mm, y=1mm, line width=0.6pt, >={Stealth[length=1.5mm]},
    cmp/.style={draw=trait, fill=white, minimum width=13mm, minimum height=4.6mm,
                inner sep=0.5mm, font=\ttfamily\bfseries\fontsize{5.8}{6.2}\selectfont},
    lbl/.style={font=\ttfamily\fontsize{5.6}{5.8}\selectfont, inner sep=0.8pt},
    titre/.style={anchor=north east, inner sep=0.8pt,
                  font=\ttfamily\bfseries\fontsize{5.8}{6.2}\selectfont},
    zone/.style={rounded corners=1mm, line width=0.5pt}]
  \draw[zone, draw=trait, fill=paper] (0,9.5) rectangle (80,34);
  \node[titre] at (79.6,34) {test = uvm\_test\_top};
  \draw[zone, draw=trait, fill=white] (2,11) rectangle (78,28.5);
  \node[titre] at (77.6,28.5) {env};
  \draw[zone, draw=cat-uvm, line width=0.8pt] (3.5,12.5) rectangle (55.5,25.5);
  \node[titre, text=cat-uvm] at (55.1,25.5) {agent};
  \node[cmp] (sqr) at (11,18) {sequencer};
  \node[cmp] (drv) at (30,18) {driver};
  \node[cmp] (mon) at (47,18) {monitor};
  \node[cmp] (scb) at (68,18) {scoreboard};
  \node[cmp, draw=cat-rnd, text=cat-rnd] (seq) at (11,31) {séquence};
  \node[cmp, minimum width=46mm] (if) at (38.5,5.5) {interface (virtual)};
  \node[cmp, minimum width=30mm, fill=paper] (dut) at (38.5,0) {DUT};
  \draw[->, cat-rnd] (seq) -- node[lbl, right, pos=0.25] {start()} (sqr);
  \draw[->] (sqr) -- node[lbl, above] {item} (drv);
  \draw[->] (drv) -- node[lbl, left, pos=0.62] {vif} (drv |- if.north);
  \draw[->] (mon |- if.north) -- node[lbl, right, pos=0.38] {vif} (mon);
  \draw[->] (mon) -- node[lbl, above] {ap} (scb);
  \foreach \x in {30,47} \draw (\x,3.2) -- (\x,2.3);
\end{schema}
\begin{optab}[19mm]
\V|uvm_object|    & items, séquences, config : créés pendant le test, sans parent \\
\V|uvm_component| & structure fixe : \V|new(name, parent)|, phases, hiérarchie
\end{optab}
\end{bloc}

\begin{bloc}{Phases}{uvm}
\begin{tab3}[27mm]{9mm}
\V|build_phase|        & func $\downarrow$ & créer les fils, lire la config \\
\V|connect_phase|      & func $\uparrow$   & relier les ports TLM \\
\V|end_of_elaboration| & func $\uparrow$   & \V|print_topology()| \\
\V|run_phase|          & task $\parallel$  & tout le temps simulé \\
\V|reset| $\cdots$ \V|main| & task         & sous-phases de run \\
\V|check report| & func $\uparrow$ & vérifier, résumer
\end{tab3}
\note{$\downarrow$ du test vers les feuilles · $\uparrow$ l'inverse ·
$\parallel$ tous les composants en parallèle.}
\sub{Squelette d'un composant, objections}
\begin{code}
class my_test extends uvm_test;
  `uvm_component_utils(my_test)     // factory
  my_env env;                       // créé en build_phase
  function new(string name, uvm_component parent);
    super.new(name, parent);        // obligatoire
  endfunction
  task run_phase(uvm_phase phase);
    wr_rd_seq seq;
    phase.raise_objection(this);    // empêche la fin
    seq = wr_rd_seq::type_id::create("seq");
    void'(seq.randomize());
    seq.start(env.agt.sqr);         // lance la séquence
    phase.drop_objection(this);     // plus aucune
  endtask                           //  -> fin du test
endclass
\end{code}
\end{bloc}

\begin{bloc}{Factory, config\_db, rapports}{uvm}
\begin{code}
drv = my_drv::type_id::create("drv", this);
it  = my_item::type_id::create("it");    // objet
// remplacer un type (dans le test, avant create)
base_drv::type_id::set_type_override(err_drv::get_type());
// set(contexte, chemin, clé, valeur) puis get
uvm_config_db#(int)::set(this, "env.agt*", "n", 4);
if (!uvm_config_db#(int)::get(this, "", "n", n))
  `uvm_fatal("CFG", "n absent")
`uvm_info("ID", $sformatf("x=%h", x), UVM_LOW)
\end{code}
\begin{optab}[17mm]
\V|`uvm_info|    & affiché si sa verbosité $\leq$ \V|+UVM_VERBOSITY| :
                   \V|UVM_NONE < LOW < MEDIUM < HIGH < FULL| \\
\V|`uvm_error|   & compté, la simulation continue (\V|`uvm_warning| idem, non bloquant) \\
\V|`uvm_fatal|   & arrête la simulation
\end{optab}
\end{bloc}
\end{colonne}\colbreak
% ======================= colonne 5 : stimulus =======================
\begin{colonne}[stimulus][col-b]
\begin{bloc}{Transaction : sequence item}{rnd}
\begin{code}
class apb_tr extends uvm_sequence_item;
  rand bit [31:0] addr, data;
  rand dir_t      dir;             // READ / WRITE
  `uvm_object_utils_begin(apb_tr)  // enregistre et
    `uvm_field_int(addr, UVM_ALL_ON | UVM_HEX)
    `uvm_field_enum(dir_t, dir, UVM_ALL_ON)
  `uvm_object_utils_end            // print, copy, compare
  function new(string name = "apb_tr");
    super.new(name);
  endfunction
endclass
\end{code}
\end{bloc}

\begin{bloc}{Séquence}{rnd}
\begin{code}
class wr_rd_seq extends uvm_sequence #(apb_tr);
  `uvm_object_utils(wr_rd_seq)
  rand int unsigned n;
  constraint c_n { n inside {[1:10]}; }
  task body();                     // le scénario
    bit [31:0] a;
    repeat (n) begin
      `uvm_do_with(req, { dir == WRITE; })
      a = req.addr;                // relit la même
      `uvm_do_with(req, { dir == READ;
                          addr == local::a; })
    end
  endtask
endclass
\end{code}
\note{\V|`uvm_do(x)| : crée, randomise, envoie l'item ou lance la sous-séquence x
(\V|start_item| + \V|randomize| + \V|finish_item|). \V|new| : comme l'item.
Lancée par le test : \V|seq.start(sqr)| (colonne 4).}
\end{bloc}

\begin{bloc}{Driver}{rnd}
\begin{code}
class apb_drv extends uvm_driver #(apb_tr);
  `uvm_component_utils(apb_drv)
  virtual apb_if vif;              // accès aux signaux
  function void build_phase(uvm_phase phase);   // + new
    super.build_phase(phase);
    if (!uvm_config_db#(virtual apb_if)::get(
          this, "", "vif", vif))   // set par tb_top
      `uvm_fatal("NOVIF", "pas de vif")
  endfunction
  task run_phase(uvm_phase phase);
    forever begin
      seq_item_port.get_next_item(req);  // bloquant
      drive(req);
      seq_item_port.item_done();         // item suivant
    end
  endtask
  task drive(apb_tr t);            // protocole APB, <=
    @(posedge vif.PCLK) vif.PSEL <= 1;  «PADDR...»
    @(posedge vif.PCLK) vif.PENABLE <= 1;
    do @(posedge vif.PCLK); while (!vif.PREADY);
    vif.PSEL <= 0;  vif.PENABLE <= 0;
  endtask
endclass
\end{code}
\end{bloc}
\end{colonne}\colbreak
% ======================= colonne 6 : observation, assemblage =======================
\begin{colonne}[observation · assemblage][col-c]
\begin{bloc}{Monitor}{obs}
\begin{code}
class apb_mon extends uvm_monitor;
  `uvm_component_utils(apb_mon)
  virtual apb_if vif;
  uvm_analysis_port #(apb_tr) ap;  // new() en build
  task run_phase(uvm_phase phase);
    apb_tr t;
    forever begin
      @(posedge vif.PCLK iff vif.PENABLE && vif.PREADY);
      t = apb_tr::type_id::create("t");
      t.addr = vif.PADDR;  «...»
      ap.write(t);                 // à tous les abonnés
    end
  endtask
endclass
\end{code}
\end{bloc}

\begin{bloc}{Scoreboard}{obs}
\begin{code}
`uvm_analysis_imp_decl(_apb)   // crée les classes
`uvm_analysis_imp_decl(_uart)  // uvm_analysis_imp_apb...
class my_scbd extends uvm_scoreboard;
  `uvm_component_utils(my_scbd)
  uvm_analysis_imp_apb  #(apb_tr,  my_scbd) apb_imp;
  uvm_analysis_imp_uart #(uart_fr, my_scbd) uart_imp;
  byte exp_q[$];                     // attendus (file)
  «new : super.new() puis apb_imp = new("apb_imp", this)»
  function void write_apb(apb_tr t);   // entrée
    if (t.dir == WRITE) exp_q.push_back(t.data[7:0]);
  endfunction
  function void write_uart(uart_fr f); // sortie
    byte e = exp_q.pop_front();
    if (f.data != e) `uvm_error("SCBD", "écart")
  endfunction
endclass
\end{code}
\end{bloc}

\begin{bloc}{Assemblage : agent, env, top}{uvm}
\begin{code}
// agent.build_phase
mon = apb_mon::type_id::create("mon", this);
if (get_is_active() == UVM_ACTIVE) begin
  sqr = apb_sqr::type_id::create("sqr", this);
  drv = apb_drv::type_id::create("drv", this);
end
// agent.connect_phase
drv.seq_item_port.connect(sqr.seq_item_export);
// env.connect_phase : monitors -> scoreboard
apb_agt.mon.ap.connect(scbd.apb_imp);
\end{code}
\sub{Module top}
\begin{code}
`include "uvm_macros.svh"
module tb_top;
  import uvm_pkg::*;  import my_pkg::*;
  logic clk = 0;  always #5 clk = ~clk;
  apb_if vif0 (clk);                // interface
  dut u_dut (.PCLK(clk), .PADDR(vif0.PADDR), «...»);
  initial begin
    uvm_config_db#(virtual apb_if)::set(null,
      "uvm_test_top.*", "vif", vif0);
    run_test();                     // +UVM_TESTNAME
  end
endmodule
\end{code}
\note{\V|vsim tb_top +UVM_TESTNAME=my_test -sv_seed random|}
\end{bloc}
\end{colonne}
\end{document}
